\clearpage
\section*{Can large condominium buildings have a floor area?}
\textit{September 23, 2026. Agent-drafted data check; keeping condominiums and reporting FAR and DUPAC as main
estimates are Jacob's decisions.}

The Assessor's condominium file covers every assessment year from 1999, but it records building and unit floor area
only for condominium buildings of fewer than 20 units. In the permit construction data,
171 of 205 condominium buildings of 20 or more units have a DUPAC and none has a FAR, while large rentals do (269 of
411, from commercial valuations). Condominiums are about two thirds of new buildings of 2 to 19 units, so dropping
them would leave a small multifamily sample of rentals; Jacob chose to keep them.

The 2022 Cook County building footprints give every building's footprint area and roof height
(\texttt{tasks/download\_building\_footprints\_2022}, 895,674 footprints over Chicago). In
\texttt{tasks/audits/footprint\_floor\_area\_check}, buildings built through 2021 were matched to the footprints
containing their parcels, and floor area was estimated as footprint area times height times a floor area per cubic
foot calibrated on rentals (0.063, about one floor per 16 feet). Against the Assessor, the estimate's median absolute
error is 13 percent for rentals under 20 units, 16 percent for condominiums under 20 units (held out of the
calibration). For rentals of 20 or more units, calibrated on the other large rentals, it is 16 percent for 20--49
units, 17 for 50--99 and 31 for 100 or more, where parking podiums, retail floors and setbacks break the link between
height and floors. Of the 160 large condominium buildings matched to footprints, 44 have 100 or more units.

Footprint floor area is not accurate enough to fill in the large condominium buildings alongside Assessor floor
area. The plan, not yet built into the analysis: FAR for buildings of fewer than 20 units, condominiums and rentals
alike, and DUPAC for all buildings, both as main estimates.

\textit{Update, same day.} Jacob chose to fill floor area for condominium buildings of 20 to 99 units rather than
drop them, marked so they can be removed. Before adopting it, the footprint estimate was examined by building height
and by footprint area relative to the lot. It fails in three identifiable situations: towers taller than 100 feet
(roof height misstates floors), buildings under 50 feet with 20 or more units (complexes spread over footprints not
matched, estimates 30 percent low), and footprints more than 1.2 times the lot (outlines that include neighbors,
estimates about 50 percent high); the first fill had put 9 W Walton at a FAR of 48. For buildings 50 to 100 feet tall
whose footprints cover half to 1.2 times their lots, calibrated on the other large rentals, the median absolute error
is 12 percent for rentals of 20 to 49 units and 8 percent for 50 to 99 units, against 31 to 34 percent for those
outside the rule; held-out condominiums under 20 units within the rule are off by 16 percent. The thresholds were
chosen on these rentals, so the held-out condominiums are the cleaner test. A summation-order error that overstated
the volume of the 84 buildings on several footprints was fixed first.

The construction data now fill floor area for 60 condominium buildings of 20 to 99 units
(\texttt{floor\_area\_source = footprint}), with FAR 1.5 to 5.8 (median 3.6, against 3.1 for Assessor-measured
buildings of that size). No other building changed; 13,499 buildings are eligible for both measures.

\textit{Large condominium buildings, researched one by one.} Of 297 condominium buildings of 20 or more units the
Assessor reports built 2006--2022, 78 had no row. Individual research (building records, permits and web sources,
\texttt{tasks/audits/construction\_hand\_checks/condominium\_site\_reviews.csv}) found 67 handled correctly: 46
represented by a later record of the same building (Trump Tower, Aqua, 1400 Museum Park), 13 completed before 2006
and 8 conversions of older buildings (Tribune Tower). The errors came from condominium buildings the Assessor
re-declares under new building numbers as units sell: early declarations counted as separate buildings (Zen Condos,
225 S Sangamon, twice) or measuring a building on a partial declaration (The Legacy at Millennium Park at 177 of 355
units; 600 N Fairbanks at 87 of 212), one row summing two towers (Lincoln Park 2550 and 2520), one separate building
dropped as a duplicate (15 N Aberdeen) and one placeholder record for a townhouse association. A general rule would
misfire where several towers share one land record, so the 21 sites were decided by hand in
\texttt{adjudication/manual\_decisions.csv}: ten rows measured on the record of the whole building, seventeen
duplicate or placeholder rows dropped, one building kept and one foundation permit joined to its building. A
suggested merge of 1411 S Michigan and 51 E 14th Street was rejected: they are separate valuations of different floor
area and land. Lux24 (24 S Morgan), a condominium building converted to rentals, stays unmeasured. Seven buildings
changed units, the file has seventeen fewer rows, and no other building changed.
